\cotefile{51}
\scanpage{13}
pas toujours ainsi, comme on voit
ainsi. On \deleted{\unsure{peut donc}} suppose $\mathrm{car.}\,k=p>0$,
on prend des ss-schémas finis $\Gamma_1,\Gamma_2$ du
schéma abélien $\deleted{A}B$, tels que $\Gamma'_1\deleted{\unsure{\Gamma_2}}\cap\Gamma'_2=(e)$,
$\Gamma_1\cap\Gamma_2\ne(e)$. On en construit facilement,
p. ex. en partant de $B$ courbe elliptique
telle que $B\otimes k$ n'ait pas de pt d'ordre
$p$ (donc $B\otimes k$ a un seul sous-groupe
fini d'ordre $p$) et $B\otimes K$ ayant deux
sous-groupes distincts d'ordre $p$ (\unsure{l'un|ceci|c'est}
\unsure{faisible|possible|faut-il|faisable} \deleted{\unsure{aux bras}} en \deleted{\unsure{égales car.}}
d'inégales caractéristiques \unsure{?|$0,p$}, \deleted{\unsure{pour que} $\unsure{k/K}$} \inserted{en prenant un}
\deleted{\unsure{pourvu} que $B\otimes K$ ait \unsure{des pts} \inserted{\unsure{de l'ordre}} \unsure{d'ordre} $p$,
on peut \unsure{en trouver} \unsure{beaucoup}, \unsure{sur le}
\unsure{deuxième cas}, \inserted{\unsure{base}} \unsure{en prenant}} \unsure{(à|à} une extension finie de $K$).
On pose $A=B/\Gamma_1\times B/\Gamma_2$, et on y envoie
$B$ diagonalement. OK.
% slanted; long arrow to the space before « Noter »
\inserted{En égales $\mathrm{car.}>0$, on prend pour $B$ un \deleted{\unsure{produit de deux} \unsure{schémas}} schéma abélien \unsure{quelc.|quel.} de dim $\geq2$, \deleted{\unsure{étant la fois}} tel que la puissance $p^{\mathrm{ième}}$ dans $\mathrm{Lie}(B)$ soit nulle, et on fait tourner un sous-espace de dim $2$ ds $\mathrm{Lie}(B)$\unsure{:} \deleted{\unsure{pour}} \inserted{\unsure{des|deux}} \unsure{droites} \unsure{faisant}, \unsure{affines|affirmer|définissant}, \unsure{partir|partant} d'une position où \unsure{ils sont distincts}.}

Noter cependant qu'en car. $0$,
\unsure{comme|pour un} \unsure{tt|$\deleted{A}$|tout} schéma en groupes de type
fini sur un corps est \unsure{séparable|appellé} \unsure{dessus|disons|donc},
% archive sheet number: 12
